% The designated English reader omits the historical front matter.
% Reconstructed and translated directly from the read-only French witness,
% PDF pp.2–11, visually checked against those pages. No page images included.
\thispagestyle{empty}
\noindent\textbf{编者}\\
Luc Illusie\\
Université de Paris-Sud\\
91405 Orsay/France
\vfill
\noindent AMS 主题分类 (1970): 14C15, 14F05, 14F20, 14G10,\\
14G13, 14H35, 14M15, 16A62, 18F30
\par\noindent\rule{\linewidth}{0.4pt}
\noindent ISBN 3-540-08248-4\quad Springer-Verlag Berlin $\cdot$ Heidelberg $\cdot$ New York\\
ISBN 0-387-08248-4\quad Springer-Verlag New York $\cdot$ Heidelberg $\cdot$ Berlin

\noindent 本书受版权保护, 保留全部权利, 无论涉及全部还是部分内容, 尤其包括翻译、重印、插图再使用、广播、影印或类似方式复制, 以及在数据库中存储的权利.

\noindent 根据德国著作权法第 54 条, 为私人使用以外的目的进行复制时, 需向出版社支付费用; 费用数额由与出版社的协议确定.

\noindent \copyright\ Springer-Verlag Berlin $\cdot$ Heidelberg 1977\\
德国印刷\\
印刷与装订: Beltz Offsetdruck, Hemsbach/Bergstr.\\
2141/3140-543210
\clearpage
\phantomsection\addcontentsline{toc}{section}{全卷导言}
\section*{导言}

本卷汇集 Grothendieck 于 1965/66 年在 IHÉS 主持的讨论班报告; 这些报告最初以油印讲义形式分发. 相对于最初版本, 重要的改动只涉及未予重印的第 II 报告, 以及完全重写并增补了编号为 III B 的附篇的第 III 报告. 除少量细节修改和脚注增补外, 其他报告均保持原样.

讨论班的核心是平展上同调中的 Lefschetz 公式 (III, III B, XII), 及其对 $L$ 函数的上同调解释的应用 (XIV). Grothendieck 在 Bourbaki 讨论班报告 [2] 中宣布的所有结果, 都在这里得到完整证明. (III, III B 和 XII) 中沿不同途径建立的迹公式, 比仅为证明 $L$ 函数的有理性所必需的形式更为一般. Deligne 的报告 (SGA $4^{1/2}$ Rapport) 中给出了一个明显更短而完整的有理性证明; 它沿用 Grothendieck 在第 XII 和 XIV 报告中的方法, 但去除了所有多余的一般性. 不过, 我们仍希望 III 和 III B 中的公式能应用于其他情形. 讨论班的其余部分包括两个关于过渡到极限的理论的报告, 从中产生 $\ell$ 进上同调 (V, VI), 以及对对偶形式体系 (I, VII) 和 Lefschetz 公式 (VII, X) 的若干补充.

下面更具体地说明本卷的内容.

第 I 报告独立于讨论班的其余部分. 它的主要结果是: 在满足某些局部附加条件的正则概形上, 若可以使用奇点消解和纯性定理, 则对与剩余特征互素的 $n$, 值为 $\mathbb Z/n\mathbb Z$ 的常值层为对偶化层 (I.3.4.1). 特别地, 借助 Hironaka 的工作和 Artin 的结果 (SGA 4 XIX), 该定理适用于特征零的卓越正则概形. 正则一维概形的情形用另一种十分初等的方法单独处理 (另见 (SGA $4^{1/2}$ Dualité)). 还应指出, Deligne 在 (SGA $4^{1/2}$ Th. Finitude) 中用既不依赖奇点消解也不依赖纯性的另一论证证明: 若 $f:X\to S$ 为维数 $0$ 或 $1$ 的正则概形上的有限型概形, 则复形 $f^!(\mathbb Z/n\mathbb Z)_S$ (其中 $n$ 与剩余特征互素) 为对偶化复形. 这个结果比第 I 报告的双对偶定理更为实用, 与后者部分相交但并不包含后者; 目前不知道它能否推广到基概形 $S$ 卓越、正则且维数大于 $1$ 的情形.

如上所述, 原题为 ``任意支撑上同调的 Künneth 公式'' 的第 II 报告未收入本卷. 该报告由 L. Illusie 根据 Grothendieck 的手稿编写, 讨论泛上同调固有性和泛局部无环性定理, 但仅在奇点消解假设下证明了这些定理. 此后 Deligne 不依赖这些假设证明了同样的结果 (SGA $4^{1/2}$ Th. Finitude), 使得这篇报告没有发表的必要. 一些补充内容, 例如关于正规交叉除子的情形, 已收入该文的附录. 再补充一点, 对集合层和非交换群层, 与该文定理类似的结果已在 Raynaud 女士的报告 (SGA 1 XIII) 中证明.

第 III 报告证明了域上有限型概形上层复形之间的上同调对应的 Lefschetz--Verdier 公式; 第 III B 报告则说明如何在若干情形下计算公式中的局部项, 例如曲线上的 Frobenius 对应. 关于这两篇报告的内容, 更详细的说明见各篇导言.

A. Grothendieck 的第 IV 报告讨论与代数闭链关联的上同调类的构造, 由 P. Deligne 编写, 已发表于 SGA $4^{1/2}$.

第 V 报告给出关于进投射系的技术性准备后, 第 VI 报告定义可构造 $\mathbb Z_\ell$-层 (分别为可构造扭常值层, 即 Deligne 在 [1] 和 (SGA $4^{1/2}$ Arcata) 中所谓的 ``光滑'' 层), 并从 $\mathbb Z/\ell^n$-层出发通过取极限, 将固有支撑高阶正像的形式体系推广到可构造 $\mathbb Z_\ell$-层. 但第 VI 报告并未定义 $\mathbb Z_\ell$-层的导出范畴, 也未构造将 $\mathbb Z/\ell^n$-层导出范畴中的运算推广而来的 $\otimes^{\mathrm L}$ 和 $\mathrm{RHom}$ 运算. 更不用说将 SGA 4 XVIII 的对偶形式体系推广到 $\mathbb Z_\ell$-层. 这个问题曾是 Jouanolou 的学位论文 (尚未发表) 的主题. 不过, 读者可以参阅 [1] 的开头, 其中定义了 $\mathbb Z_\ell$-层的 $\Ext^1$.

第 VII 报告独立于讨论班的其余部分, 无疑是其中最有意思的报告之一. 它包括 a) 若干标准代数簇 (仿射空间、射影空间、旗簇) 的上同调计算, b) 平展上同调中的 Chern 类理论, c) 自交公式 $i^*i_*x=x c_d(N)$ 在平展上同调中的证明 (分别在 Chow 环中的证明; 后一情形的证明出自 Mumford), 以及该公式的若干应用 (爆破代数簇的上同调计算、Gauss--Bonnet 公式).

第 VII 报告包含导出范畴中的有限性条件和 Grothendieck 群的一系列基本结论. 这些问题在 (SGA 6 I, IV) 中于更一般的框架下重新讨论.
% The preceding VII is VII in the French Introduction (p.V), although VIII is intended.

J.-P. Serre 的第 IX 报告题为 ``Brauer 理论导引'', 已发表于 [4], 本卷不再重印. 它的一个主要结果 (出自 Swan) 是存在一个以 ``Swan 特征标'' 为特征标的投射模; 第 X 报告使用了这个结果, 并证明曲线上的层的一个 Euler--Poincaré 公式, 比 Raynaud 的报告 [3] 中的公式略为一般.

I. Bucur 编写的第 XI 报告已不复存在, 因为稿件在一次搬家中遗失, 作者也没有副本. 该报告讨论 Grothendieck 的交换迹理论, 推广了 Stallings [5] 的理论. (III B 5, 6) 在层论框架下给出了这一理论的更精细版本. 第 XII 和 III B 报告给出了对 Lefschetz 公式的应用. 第 XII 报告的迹公式独立于第 III 报告的一般公式而证明, 但 (III B 6) 说明其中的局部项确实是一般公式中的局部项, 且一般公式蕴含该公式.
% French p.VI says traces commutatives; source wording retained, despite subject.

第 XIII 报告的缺席仅仅由于编号问题: 最初计划把第 X 至 XII 报告的内容分布在四篇报告中.

第 XIV 报告 (有时也编号为 XV) 介绍 Frobenius 对应的一般性质, 特别是 ``几何'' Frobenius 与 ``算术'' Frobenius 的比较, 并从第 XII 报告 (或 III B) 的迹公式出发, 给出 Grothendieck 对 $L$ 函数有理性的证明. 其中两个要点是归结为曲线情形, 以及通过取极限归结为关于 $\mathbb Z/\ell^n$-层的公式.

在口头讨论班中, Grothendieck 曾报告尚未解决的问题并提出若干猜想. 这篇本应结束讨论班的报告未能写成, 令人遗憾. 他那篇十分出色的开场报告也未写成; 那篇报告综述了不同背景 (拓扑、复解析、代数) 中的 Euler--Poincaré 公式和 Lefschetz 公式.

我感谢 P. Deligne 说服我在第 III 报告的新版本中写出 Lefschetz--Verdier 公式的证明, 从而消除了出版本讨论班讲义的一项障碍. 对于他建议的写作改进, 以及他在准备本卷交付出版社过程中给予我的帮助, 我深表感谢. 我也感谢 Lièvremont 女士, 她在很短时间内细心完成了第 III 和 III B 报告的打字工作.

结束这篇导言之前, 我必须向 I. Bucur 表达悼念: 他于 1976 年 9 月因癌症逝世. 那些像我一样有幸以他为友的人, 不会忘记他细致的体贴、带着微笑的幽默, 以及他谈吐的魅力.

\begin{flushright}
Paris, 1977 年 2 月 19 日\\[1em]
Luc Illusie
\end{flushright}

\subsection*{参考文献}
\begin{enumerate}[label={[\arabic*]}]
\item Deligne, P., \emph{La conjecture de Weil II}, 待刊.
\item Grothendieck, A., \emph{Formule de Lefschetz et rationalité des fonctions $L$}, Séminaire Bourbaki 64/65, no.~279, 载于 \emph{Dix exposés sur la cohomologie des schémas}, North Holland, Pub. Co., 1968.
\item Raynaud, M., \emph{Caractéristique d'Euler--Poincaré d'un faisceau et cohomologie des variétés abéliennes (d'après Ogg--Shafarevitch et Grothendieck)}, Séminaire Bourbaki 64/65, no.~286, 载于 \emph{Dix exposés sur la cohomologie des schémas}, North Holland Pub. Co., 1968.
\item Serre, J.-P., \emph{Représentations linéaires des groupes finis}, Hermann, 重编第 2 版, 1971.
\item Stallings, J., \emph{Centerless groups -- an algebraic formulation of Gottlieb's theorem}, Topology, 4, 第 129--134 页, 1965.
\end{enumerate}

\clearpage
\phantomsection\addcontentsline{toc}{section}{SGA 书目}
\section*{SGA 书目}
SGA = Séminaire de Géométrie Algébrique du Bois-Marie (博瓦马里代数几何讨论班).
\begin{description}[style=nextline]
\item[SGA 1] 平展覆盖与基本群, A. Grothendieck 著, Lecture Notes in Mathematics no.~224, Springer-Verlag, 1971.
\item[SGA 2] 凝聚层的局部上同调与局部和整体 Lefschetz 定理, A. Grothendieck 著, North Holland Pub. Co., 1968.
\item[SGA 3] 群概形 I, II, III, M. Demazure 与 A. Grothendieck 著, Lecture Notes in Mathematics nos.~151, 152, 153, Springer-Verlag, 1970.
\item[SGA 4] 拓扑斯理论与概形的平展上同调, M. Artin, A. Grothendieck, J.-L. Verdier 著, Lecture Notes in Mathematics nos.~269, 270, 305, Springer-Verlag, 1972--73.
\item[SGA $4^{1/2}$] 平展上同调, P. Deligne 著, Lecture Notes in Mathematics no.~569, Springer-Verlag, 1977.
\item[SGA 5] $\ell$ 进上同调与 $L$ 函数, A. Grothendieck 著, Lecture Notes in Mathematics no.~589, Springer-Verlag, 1977.
\item[SGA 6] 相交理论与 Riemann--Roch 定理, P. Berthelot, A. Grothendieck, L. Illusie 著, Lecture Notes in Mathematics no.~225, Springer-Verlag, 1971.
\item[SGA 7] 代数几何中的单值群, I 由 A. Grothendieck 著, II 由 P. Deligne 与 N. Katz 著, Lecture Notes in Mathematics nos.~288, 340, Springer-Verlag, 1972, 1973.
\end{description}
\clearpage
\section*{原版目录\footnote{关于第 II, IV, IX, XI, XIII 报告缺席的说明, 见导言.}}
% Historical source page numbers retained separately from generated Chinese ToC.
\begin{description}[style=nextline]
\item[第 I 报告\hfill 1] 对偶化复形, A. Grothendieck 著, L. Illusie 编写.
\item[第 III 报告\hfill 73] Lefschetz 公式, A. Grothendieck 著, L. Illusie 编写.
\item[第 III B 报告\hfill 138] 局部项的计算, L. Illusie 著.
\item[第 V 报告\hfill 204] $J$ 进投射系, J.-P. Jouanolou 著.
\item[第 VI 报告\hfill 251] $\ell$ 进上同调, J.-P. Jouanolou 著.
\item[第 VII 报告\hfill 282] 若干经典概形的上同调与 Chern 类的上同调理论, J.-P. Jouanolou 著.
\item[第 VIII 报告\hfill 351] 阿贝尔范畴与三角范畴的类群, 完美复形, A. Grothendieck 著, I. Bucur 编写.
\item[第 X 报告\hfill 372] 平展上同调中的 Euler--Poincaré 公式, A. Grothendieck 著, I. Bucur 编写.
\item[第 XII 报告\hfill 407] 代数几何中的 Nielsen--Wecken 公式与 Lefschetz 公式, A. Grothendieck 著, I. Bucur 编写.
\item[第 XIV = XV 报告\hfill 442] Frobenius 态射与 $L$ 函数的有理性, C. Houzel 著.
\item[术语索引\hfill 481]\mbox{}
\item[符号索引\hfill 483]\mbox{}
\end{description}
\begin{center}
\begin{tikzpicture}[>=stealth]
\node (independent) at (0,0) {I, VII, VIII};\node (x) at (3,0) {X};\draw[->] (independent)--(x);
\node (iii) at (5,1) {III};\node (iiib) at (7.5,1) {III B};\draw[->](iii)--(iiib);
\node (v) at (5,0) {V};\node (vi) at (7.5,0) {VI};\draw[->](v)--(vi);
\node (iiib5) at (5,-1) {III B 5};\node (xii) at (7.5,-1) {XII};\draw[->](iiib5)--(xii);
\node (xiv) at (10,0) {XIV};\draw[->](iiib)--(xiv);\draw[->](vi)--(xiv);\draw[->](xii)--(xiv);
\end{tikzpicture}
\end{center}
\clearpage
